\documentclass[11pt]{article}

% ----- Packages -----
\usepackage[margin=1in]{geometry}
\usepackage{amsmath, amssymb, amsthm}
\usepackage{braket}          % \ket, \bra, \braket, \ketbra
\usepackage{mathtools}
\usepackage{bm}
\usepackage{enumitem}
\usepackage{xcolor}
\usepackage{hyperref}
\hypersetup{colorlinks=true, linkcolor=blue!55!black, urlcolor=blue!55!black}

% ----- Convenience macros -----
\newcommand{\Tr}{\operatorname{Tr}}
\newcommand{\onenorm}[1]{\left\lVert #1 \right\rVert_{1}}
\newcommand{\R}{\mathbb{R}}
\newcommand{\C}{\mathbb{C}}

\title{\textbf{Quantum Information Theory}\\[4pt]
Fidelity of Quantum States:\\
An Introduction to Uhlmann's Theorem}
\author{}
\date{Talk 3 --- July 9, 2026}

\begin{document}
\maketitle

%===============================================================
\section*{Review}
%===============================================================

\textbf{Classical discrimination} of a classical random variable $x$ drawn
from one of the pdf's $p(x)$ or $q(x)$.
Rule: given observation $x$, choose the greater of
\[
\begin{cases}
p(x) \longrightarrow x\sim p,\\[2pt]
q(x) \longrightarrow x\sim q,
\end{cases}
\qquad
\Big(\text{gives maximum prob.\ of success}\Big).
\]

\textbf{Quantum discrimination} between two mixed states $\rho,\sigma$: we
designed the optimal POVM $\{E,\,I-E\}$. The probability of successful guessing
in both cases is
\[
P^{\mathrm{succ}}_{\mathrm{classical}}
= \frac12\!\left(1 + \frac12\onenorm{p-q}\right),
\qquad
P^{\mathrm{succ}}_{\mathrm{quantum}}
= \frac12\!\left(1 + \frac12\onenorm{\rho-\sigma}\right),
\]
where
\[
\onenorm{A} = \sum_i |\lambda_i(A)| = \Tr\sqrt{A^{\dagger}A}.
\]
We defined the \emph{trace distance} between states $\rho$ and $\sigma$,
\[
T(\rho,\sigma) = \frac12\onenorm{\rho-\sigma},
\]
which is an operational definition in terms of a classification task, and also
a nice mathematical definition with an interpretation as a proper metric. Other
operational definitions are possible and useful in their own particular
applications.

\medskip
\noindent\textbf{Today we look at fidelity, $F(\rho,\sigma)$, between two
states.}

%===============================================================
\section{Fidelity as a Generalization of Overlap}
%===============================================================

Fidelity is a generalization of the overlap between two pure states,
$\braket{\psi|\phi}$; we want something similar for mixed states.

\paragraph{Claim.}
For pure states there is a direct relation between overlap and
distinguishability:
\[
\braket{\psi|\phi}\ \sim\
\frac12\onenorm{\,\ket{\psi}\!\bra{\psi} - \ket{\phi}\!\bra{\phi}\,}
\qquad(\text{not equal!}).
\]

\paragraph{Why.}
Take the $2$-dimensional subspace spanned by $\ket{\psi}$ and $\ket{\phi}$.
Take basis vector $\ket{e_1}\equiv\ket{\psi}$ and any orthonormal $\ket{e_2}$.
Then
\[
\ket{\psi} = \begin{pmatrix}1\\0\end{pmatrix},
\qquad
\ket{\phi} = \cos\theta\,\ket{e_1} + \sin\theta\,\ket{e_2}
= \begin{pmatrix}\cos\theta\\ \sin\theta\end{pmatrix},
\quad\text{where } \braket{\psi|\phi} = \cos\theta.
\]
So, writing $c=\cos\theta$ and $s=\sin\theta$,
\begin{align*}
\ket{\psi}\!\bra{\psi} - \ket{\phi}\!\bra{\phi}
&= \begin{pmatrix}1\\0\end{pmatrix}\!\begin{pmatrix}1&0\end{pmatrix}
- \begin{pmatrix}c\\s\end{pmatrix}\!\begin{pmatrix}c&s\end{pmatrix}\\[4pt]
&= \begin{pmatrix}1&0\\0&0\end{pmatrix}
- \begin{pmatrix}c^{2}&cs\\ cs&s^{2}\end{pmatrix}
= \begin{pmatrix}s^{2}&-cs\\ -cs&-s^{2}\end{pmatrix},
\end{align*}
and we need $\sum_{i=1}^{2}|\lambda_i|$ of its eigenvalues.

Recall the characteristic equation for a $2\times 2$ matrix,
$\lambda^{2} - \Tr(A)\,\lambda + \det(A) = 0$. Here $\Tr = 0$ and
\[
\det = -s^{4} - c^{2}s^{2} = -s^{2}(s^{2}+c^{2}) = -\sin^{2}\theta,
\]
so
\[
\lambda^{2} - \sin^{2}\theta = 0
\quad\Longrightarrow\quad
\lambda = \pm\sin\theta.
\]
Therefore $|\lambda_1| + |\lambda_2| = 2\sin\theta$, and
\[
\boxed{\;
\frac12\onenorm{\,\ket{\psi}\!\bra{\psi} - \ket{\phi}\!\bra{\phi}\,}
= \sin\theta = T(\ket{\psi},\ket{\phi}),
\qquad
\text{but}\quad |\!\braket{\psi|\phi}\!|^{2} = \cos^{2}\theta.
\;}
\]
This is a simple $1{:}1$ relation for pure states. This
\emph{squared-overlap} is what we call fidelity for pure states.

%===============================================================
\section{General Definition of Fidelity}
%===============================================================

For two mixed states $\rho,\sigma$ in general,
\[
\boxed{\; F(\rho,\sigma) = \onenorm{\sqrt{\rho}\,\sqrt{\sigma}}^{\,2}. \;}
\]

%---------------------------------------------------------------
\subsection{Recovering the pure-state definition}
%---------------------------------------------------------------

\paragraph{Claim.}
This recovers the pure-state definition above: $\cos^{2}\theta$.

\paragraph{Proof.}
Take $\rho=\ket{\psi}\!\bra{\psi}$ and $\sigma=\ket{\phi}\!\bra{\phi}$; we want
to show $\onenorm{\sqrt{\rho}\,\sqrt{\sigma}}^{\,2} = |\!\braket{\psi|\phi}\!|^{2}$.
Note that for a rank-arbitrary projection such as $\Pi=\ket{\psi}\!\bra{\psi}$
we have $\Pi^{2}=\Pi$, hence $\Pi=\sqrt{\Pi}$. So
\begin{align*}
F(\psi,\phi)
&= \onenorm{\sqrt{\ket{\psi}\!\bra{\psi}}\,\sqrt{\ket{\phi}\!\bra{\phi}}}^{\,2}
= \onenorm{\,\ket{\psi}\!\bra{\psi}\,\ket{\phi}\!\bra{\phi}\,}^{\,2}\\[4pt]
&= \onenorm{\,\underbrace{\braket{\psi|\phi}}_{\in\,\C}\,\ket{\psi}\!\bra{\phi}\,}^{\,2}
= |\!\braket{\psi|\phi}\!|^{2}\,\onenorm{\,\ket{\psi}\!\bra{\phi}\,}^{\,2},
\end{align*}
using $\onenorm{cA} = |c|\,\onenorm{A}$. It remains to show the last one-norm
equals $1$.

\medskip
\noindent\emph{One-norm of the rank-one operator $\ket{\psi}\!\bra{\phi}$.}
Take $X=\ket{\psi}\!\bra{\phi}$, so $\onenorm{X}=\Tr\sqrt{X^{\dagger}X}$ from
the definition. Then
\[
X^{\dagger}X = \ket{\phi}\!\underbrace{\braket{\psi|\psi}}_{=\,1}\!\bra{\phi}
= \ket{\phi}\!\bra{\phi},
\]
and since $\ket{\phi}\!\bra{\phi}$ is a rank-one projection,
$\sqrt{X^{\dagger}X}=\ket{\phi}\!\bra{\phi}$. Therefore
\[
\Tr\sqrt{X^{\dagger}X} = \Tr\ket{\phi}\!\bra{\phi} = \Tr\braket{\phi|\phi}
= \Tr 1 = 1
\quad\Longrightarrow\quad
\onenorm{\,\ket{\psi}\!\bra{\phi}\,} = 1.
\]
Hence
\[
\boxed{\; F(\ket{\psi},\ket{\phi}) = |\!\braket{\psi|\phi}\!|^{2}, \;}
\]
as expected. \hfill$\square$

%---------------------------------------------------------------
\subsection{Fidelity of a pure and a mixed state}
%---------------------------------------------------------------

\paragraph{Claim.}
$F(\ket{\psi},\rho) = \braket{\psi|\rho|\psi}$.

\paragraph{Proof.}
With $\Pi=\ket{\psi}\!\bra{\psi}$,
\[
F(\psi,\rho)
= \onenorm{\sqrt{\ket{\psi}\!\bra{\psi}}\,\sqrt{\rho}}^{\,2}
= \onenorm{\,\ket{\psi}\!\bra{\psi}\,\sqrt{\rho}\,}^{\,2}
= \onenorm{\Pi\sqrt{\rho}}^{\,2}
\qquad(\Pi\text{ a projection}).
\]
Using $\onenorm{A}=\Tr\sqrt{AA^{\dagger}}$,
\[
\onenorm{\Pi\sqrt{\rho}}
= \Tr\sqrt{\Pi\sqrt{\rho}\,\sqrt{\rho}\,\Pi}
= \Tr\sqrt{\Pi\rho\Pi}.
\]
Now
\[
\Pi\rho\Pi = \ket{\psi}\!\bra{\psi}\rho\ket{\psi}\!\bra{\psi}
= \underbrace{\braket{\psi|\rho|\psi}}_{\in\,\C}\,\ket{\psi}\!\bra{\psi},
\]
so
\[
\Tr\sqrt{\Pi\rho\Pi}
= \sqrt{\braket{\psi|\rho|\psi}}\;\Tr\sqrt{\ket{\psi}\!\bra{\psi}}
= \sqrt{\braket{\psi|\rho|\psi}}\;\underbrace{\Tr\ket{\psi}\!\bra{\psi}}_{=\,1}
= \sqrt{\braket{\psi|\rho|\psi}}.
\]
Thus $\onenorm{\Pi\sqrt{\rho}} = \sqrt{\braket{\psi|\rho|\psi}}$, and
\[
\boxed{\; F(\psi,\rho) = \braket{\psi|\rho|\psi}. \;}
\]
\hfill$\square$

%---------------------------------------------------------------
\subsection{Operational interpretation}
%---------------------------------------------------------------

Both cases have a similar operational interpretation as the probability of a
measurement outcome for the POVM $\{\ket{\psi}\!\bra{\psi},\; I-\ket{\psi}\!\bra{\psi}\}$:
\[
\begin{aligned}
F(\psi,\phi) &= |\!\braket{\psi|\phi}\!|^{2}
&&= \text{prob.\ of outcome } \psi \text{ when given state } \ket{\phi};\\
F(\psi,\rho) &= \braket{\psi|\rho|\psi}
&&= \text{prob.\ of outcome } \psi \text{ when given state } \rho.
\end{aligned}
\]

%===============================================================
\section{Mixed States}
%===============================================================

For two mixed states,
\[
F(\rho,\sigma) = \onenorm{\sqrt{\rho}\,\sqrt{\sigma}}^{\,2}.
\]

%---------------------------------------------------------------
\subsection{Commuting operators}
%---------------------------------------------------------------

For commuting operators it is easy to show
\[
F(\rho,\sigma) = \left(\sum_i \sqrt{p_i\,q_i}\right)^{2},
\]
where $p_i$ are the eigenvalues of $\rho$ and $q_i$ are the eigenvalues of
$\sigma$. This is the Bhattacharyya distance of the classical probability
distributions $\{p_i\},\{q_i\}$.

%---------------------------------------------------------------
\subsection{Non-commuting mixed states}
%---------------------------------------------------------------

How do we measure an ``overlap'' of mixed states? Recall the
\emph{purification} of mixed states. Take a mixed state $\rho_A$ on a system
$A$; add a reference system $R$ and produce a \emph{pure state}
$\ket{\psi}_{RA}$ on the joint system $RA$ such that discarding (tracing out)
$R$ on the joint state produces $\rho_A$:
\[
\rho_A \xrightarrow{\ \text{purification}\ } \ket{\psi}_{RA}
\qquad\text{s.t.}\qquad
\Tr_R \ket{\psi}_{RA}\!\bra{\psi}_{RA} = \rho_A.
\]
Do this purification on both $\rho$ and $\sigma$, giving
$\rho\to\ket{\psi_\rho}$ and $\sigma\to\ket{\psi_\sigma}$. Try
\[
F(\sigma_A,\rho_A) \overset{?}{=} |\!\braket{\psi_\rho|\psi_\sigma}\!|^{2}.
\]
Not quite! It turns out this is \emph{not unique} --- but there \emph{is} a
unique maximum. So
\[
\boxed{\;
F(\rho_A,\sigma_A) = \onenorm{\sqrt{\rho}\,\sqrt{\sigma}}^{\,2}
\quad\text{(by definition)}
\qquad\text{ALSO}\qquad
= \max_{\text{purifications}} |\!\braket{\psi_\rho|\psi_\sigma}\!|^{2}.
\;}
\]
\textbf{This is Uhlmann's Theorem.}

%===============================================================
\section{Review of Purification}
%===============================================================

Recall: (1) a pure state $\ket{\psi}$; and (2) a mixed state $\rho$ with
$\rho\ge 0$ and $\Tr\rho=1$.

We can view a mixed state $\rho_A$ on a Hilbert space $A$ as arising from a
pure state $\ket{\psi_\rho}_{RA}$ on an extended Hilbert space $R\otimes A$,
\[
\rho_{RA} = \ket{\psi_\rho}_{RA}\!\bra{\psi_\rho}_{RA},
\]
then discarding system $R$ and keeping only $A$. Here $\ket{\psi}_{RA}$ is the
\emph{purification} of $\rho_A$. This is \emph{not unique}: infinitely many
different $\ket{\psi_\rho}_{RA}$ give the same $\rho_A$ after discarding $R$.

%---------------------------------------------------------------
\subsection{Construction}
%---------------------------------------------------------------

Given $\rho_A$, take its eigendecomposition
$\rho_A = \sum_i p_i\,\ket{e_i}\!\bra{e_i}$, and define
\[
\boxed{\;
\ket{\psi_\rho}_{RA} = \sum_i \sqrt{p_i}\,\ket{i}_R\,\ket{e_i}_A,
\;}
\]
where $\ket{i}_R$ is an arbitrary orthonormal basis in another Hilbert space
with $\dim(R)=\dim(A)$.

\paragraph{Check.}
Trace out system $R$ by applying $\Tr_R \otimes I_A$ (trace on $R$, identity on
$A$) to $\ket{\psi_\rho}_{RA}\!\bra{\psi_\rho}_{RA}$:
\begin{align*}
\Tr_R \sum_{i,j} \sqrt{p_i p_j}\,\ket{i}\!\bra{j}_R \otimes \ket{e_i}\!\bra{e_j}
&= \sum_{i,j} \sqrt{p_i p_j}\,
\underbrace{\Tr\!\big(\ket{i}\!\bra{j}_R\big)}_{=\,\delta_{ij}}
\otimes \ket{e_i}\!\bra{e_j}\\
&= \sum_i p_i\,\ket{e_i}\!\bra{e_i} = \rho_A. \qquad\checkmark
\end{align*}

%---------------------------------------------------------------
\subsection{Non-uniqueness and Uhlmann's theorem}
%---------------------------------------------------------------

Purifications are \emph{not unique}. In
\[
\ket{\psi_\rho}_{RA} = \sum_i \sqrt{p_i}\,\ket{i}_R \otimes \ket{e_i}_A,
\]
the $\ket{e_i}$ form an ONB on $A$ (eigenvectors of $\rho_A$) while $\ket{i}_R$
is an arbitrary ONB on another Hilbert space. We are free to rotate the ONB on
$R$ with a unitary $U_R$: that is, $(U_R\otimes I_A)\ket{\psi_\rho}_{RA}$ is
also a purification for $\rho_A$, given a purification $\ket{\psi_\rho}_{RA}$.

\paragraph{Relation to fidelity between states $\rho$ and $\sigma$.}
Given purifications $\rho_A\to\ket{\psi_\rho}_{RA}$ and
$\sigma_A\to\ket{\psi_\sigma}_{RA}$, we can define \emph{a} fidelity between
$\rho$ and $\sigma$ as the overlap of the purifications,
$|\!\braket{\psi_\rho|\psi_\sigma}\!|^{2}$. Define \emph{the} fidelity between
$\rho$ and $\sigma$ as
\[
\boxed{\;
F(\rho,\sigma) = \max_{\text{purifications}}
|\!\braket{\psi_\rho|\psi_\sigma}\!|^{2}.
\;}
\]
One then shows this also equals $\onenorm{\sqrt{\rho}\,\sqrt{\sigma}}^{\,2}$ ---
with no optimization procedure required, via a direct calculation.

\medskip
\noindent\textbf{This is Uhlmann's Theorem.}

\end{document}
